\renewcommand{\abstractname}{ÖZ}
\begin{abstract} 
\addcontentsline{toc}{section}{ÖZ} 
\begin{center}
\textbf{TÜRKİYE'DE İŞ GÜCÜ PİYASASI: POLARİZASYONUN
AMPİRİK ARAŞTIRMASI}
\par\end{center}

\begin{center}
\textbf{Evren Gülser}
\par\end{center}

\begin{center}
\textbf{Ağustos, 2023}
\par\end{center}

Çalışma Ekonomisi alanında, istihdamdaki değişimlerin
zaman içindeki yapısını inceleyen bir literatür giderek büyüyor.
Arz yönlü bakış açısı, iş gücünün yaşlanması,
kadınların iş gücüne artan katılım oranı ve eğitimde
geçirilen sürenin artmasına bağlı açıklamaları önemli
buluyor. Talep yönlü perspektif ise teknolojik ilerleme ve küresel
ticaretin genişlemesini temel alıyor. 2000'li yılların başlarında,
istihdamdaki değişimin genel özelliklerini tanımlamak için
medyan ücret ile belirlenen yetenek dağılımının kuyruklarında
belirgin istihdam artışına işaret eden ve \textquotedblleft iş
kutuplaşması\textquotedblright{} adı verilen bir olgu ortaya
çıktı. Teknolojik ilerlemenin, değişen iş gücünün
yapısındaki önemli etkisi üzerinde bir uzlaşı olmasına
rağmen, bu etkinin dinamikleri ve arz yönlü etkilerin yüksek katkı
düzeyine işaret eden çalışmalar, araştırmacıları
çeşitli emek piyasalarına ait verileri analiz etmeye yöneltti
ve farklı emek piyasalarından gelen kanıtların önemini artırdı.
Bu çalışma, Türkiye emek piyasasına ait kanıtlarla bu alana
katkı yapmayı amaçlamaktadır. Çalışmada temel veri seti
olarak TÜİK tarafından sağlanan İş gücü İstatistikleri
yıllık kesitler arası mikro veri setini kullanılır ve verinin
getirdiği kısıtlar dahilinde çalışma periyodu 2004-2022
arası olarak belirlenmiştir. Veri seti, 2004-2012 ve 2012-2022
olmak üzere farklı özellikler gösteren iki parçaya bölünüp havuzlanmıştır.
Literatürle uyumlu olarak göreve dayalı bir analizi kolaylaştırmak
için yeni veri setleri oluşturmuştur. Hazırlanan veri setlerinden
birisi, çalışanların işlerindeki görevleri ile ilgili
bireysel değerlendirmelerini kapsarken diğerleri görevlerin
zaman içinde değişen yapısını incelemeye olanak vermektedir. 

İlk dönemde, beceri düzeyi (medyan ücretle temsil edilen) ile
istihdam artışı arasında neredeyse pozitif doğrusal
bir ilişki bulunmaktadır. Buna karşılık, önemli reel
ücret artışlarının potansiyel olarak ikinci dönemde emek
piyasasında piyasa güçlerinin ilişkisini etkilemesine rağmen,
mesleklerin iş gücü içindeki paylarının değişimi kutuplaşma
formuna daha yakın bir görüntü çiziyor. Kanıtlara göre, değişen
sektörel kompozisyon, istihdam kaymalarının şekillenmesinde
rol oynuyor. Elde edilen bulgular, 2004-2012 yılları arasında
ücret polarizasyonu, 20012-2022 yılları arasındaki dönemde
de iş polarizasyonu olduğuna işaret etmektedir. Mesleklere
özgü özellikler, ilk dönemdeki ücret polarizasyonunu açıklama konusunda
yeterlidir. Meslek ve işçi düzeyinde görev ölçümlerinin, ücret
tahmini denklemleri için anlamlı indikatörler olduğu görülmüştür.

\end{abstract}

\textbf{Anahtar Kelimeler: }iş gücü dinamikleri, iş kutuplaşması,
ücret kutuplaşması, ikame esnekliği, görev tabanlı analiz,
emek piyasası, emek talebi, sektörel kompozisyon, yetenek Dağılımı

\textbf{JEL kodları: }J21, J24, J31
\begin{center}
\pagebreak
\par\end{center}
